\chapter{Background and Literature Review} \label{chapter:background}

\section{Material Interfaces} \label{section:material_interfaces}

\label{para:interface_section_intro} % ---------- %

Building on our foundational definitions, we now explore the key functional regions of a material interface, the
engineering properties that govern strength, adhesion, and stability, and the principal physical effects arising at
these junctions, electronic phenomena and energy-transport processes.

We then discuss design parameters for tailoring interfacial behaviour, the role of crystallographic orientation and
polymorph-driven reconstruction, and the generation and prediction of interfacial defects.

Finally, we highlight representative applications, from gate dielectric layers and transition metal dichalcogenides
(TMDC) heterojunctions to energy-storage interfaces, thermoelectric contacts, and optoelectronic devices, and conclude
with current active modelling challenges.

\label{para:interfacial_engineering_properties} % ---------- %

In materials engineering, interfaces govern a wide array of properties including adhesion, corrosion resistance,
mechanical strength, and diffusion kinetics.

The strength and ductility of composite systems or coatings often depend on interfacial bonding and lattice
compatibility.

Tailoring interfacial chemistry or morphology, through methods such as surface functionalisation, alloying, or atomic
layer deposition, is often essential for achieving desired macroscopic performance metrics.

\label{para:electronic_interface_effects} % ---------- %

In electronic devices, interfaces critically influence charge transport, band alignment, and contact resistance.

Semiconductor heterojunctions, metal-semiconductor contacts, and dielectric boundaries are fundamental to the operation
of transistors, diodes, and capacitors.

Localised changes in atomic registry, roughness, or defect density at the interface can significantly alter tunnelling
behaviour, barrier heights, and mobility.

Where interfacial dipoles, charge redistribution, and electrostatic reconstruction result in complex electronic
responses not captured by bulk models.

\label{para:energy_interfaces} % ---------- %

Energy systems are similarly interface-dominated.

In lithium-ion batteries, for example, the solid-electrolyte interphase (SEI) acts as a metastable interfacial layer
that dictates both stability and charge transport.

The performance and longevity of solid-state batteries, fuel cells, and thermoelectric devices can depend on the
engineered properties of interfaces, which affect ionic conductivity, catalytic activity, thermal transport, and
electronic insulation.

In photovoltaics, interfacial band offsets between donor and acceptor layers, or between absorber materials and charge
transport layers, govern charge separation efficiency and open-circuit voltage.

\label{para:interface_design_parameters} % ---------- %

Overall, interfaces are not incidental features but active design parameters.

They serve as sites of emergent phenomena, such as interfacial polarisation, bandgap renormalisation, and phonon
scattering, that cannot be predicted from bulk properties alone.

As such, mastering interfacial science is a gateway to developing high-performance materials and devices.

\label{para:interface_definition} % ---------- %

The structural features that influence interface behaviour include crystallographic orientation, the nature of the phase
boundary, and the presence of defects.

\label{para:orientation_effects} % ---------- %

The relative orientation of crystals on either side of an interface can significantly affect the atomic coordination,
bond continuity, and symmetry properties across the interface.

Epitaxial alignment, where lattice planes across the interface maintain a well-defined relationship, typically results
in low-strain, coherent interfaces.

In contrast, a mismatch in orientation can lead to misfit dislocations or interfacial strain, influencing charge carrier
transport, interfacial states, and mechanical stability.

Orientation also affects the formation energy and electronic reconstruction at the interface, which can lead to
phenomena like interface-induced states or altered band alignments.

\label{para:polymorph_interface_reconstruction} % ---------- %

When the interface forms between distinct structural phases, such as cubic and hexagonal polymorphs, unique bonding
geometries and strain fields arise.

For instance, in layered 2D|3D systems like \ce{HfS2}|\ce{HfO2}, the band alignment changes depending on whether the
materials are stacked or laterally connected, with lateral interfaces showing increased reconstruction due to bonding
mismatch between the different crystalline motifs.

These reconstructions are not merely geometric; they result in different electrostatic and electronic behaviours,
influencing band offsets and charge transfer dynamics.

\label{para:interfacial_defects} % ---------- %

Defects, including vacancies, interstitials, dislocations, and grain boundaries, are often intrinsic to interface
formation and can profoundly modify interfacial properties.

They may serve as scattering centres, recombination sites, or sources of local doping.

In grain boundaries of polycrystalline materials, for example, variations in local stoichiometry and strain can create
regions of localised metallicity or enhanced dielectric response, as shown in systems exhibiting colossal permittivity.

In some systems, such as CCTO (\ce{CaCu3Ti4O12}), interfacial defects and microstructural features have been shown to
dominate the macroscopic dielectric properties.

\label{para:defect_ml_prediction} % ---------- %

In computational studies, predicting the impact of such defects requires careful treatment of their spatial distribution,
charge state, and interaction with surrounding lattice atoms.

Modern approaches like the RAFFLE algorithm allow for machine learning-driven structure prediction that accounts for
void-filling and local bonding environments, enabling more realistic models of defect-laden interfaces.

\label{para:interface_applications} % ---------- %

Depending on their nature, interfaces may enhance or impede material performance.

This subsection surveys several application domains in which understanding and engineering interfacial properties is
essential.

\label{para:gate_dielectric_interface} % ---------- %

In semiconductor transistor devices, the interface between the gate dielectric and the semiconducting channel is a
critical design parameter.

% ---------- %

As traditional \ce{SiO2} dielectrics are replaced by high-$\kappa$ materials such as \ce{HfO2} to meet miniaturisation
demands, understanding charge transfer and band alignment across interfaces like \ce{HfS2}| \ce{HfO2} becomes essential
for ensuring device reliability and performance.

Notably, in such 2D|3D heterostructures, the type and strength of bonding (e.g. van der Waals vs. chemical) directly
influences charge redistribution and can cause significant deviation from band alignments predicted by Anderson's rule,
the vacuum levels of the two materials must be aligned.

\label{para:tmdc_interface_band_alignment} % ---------- %

In layered 2D electronics, especially those using TMDCs, the stacking geometry, layer
orientation, and interfacial disorder all strongly influence the band alignment.

These factors affect carrier dynamics, optical transitions, and switching behaviour.

Traditional models such as Anderson's rule have proven insufficient; corrections using terms such as $\Delta E_\Gamma$ and
$\Delta E_{\text{IF}}$ have been developed to capture hybridisation and interface dipole contributions more accurately.

% todo: add formula for detla E_\Gamma and delta E_IF

These corrections allow for predictive modelling of heterostructure behaviour and support targeted band offset
engineering.

\label{para:energy_storage_interface} % ---------- %

Interfaces also govern performance in energy storage systems.

In all-solid-state batteries, for example, the interface between electrode and solid electrolyte dictates ion mobility,
chemical stability, and space charge layer formation.

Such interfacial processes can lead to resistive interphases or dendrite formation, reducing battery life and safety.

Density Functional Theory (DFT) modelling has shown that interfacial reconstructions and electrostatic gradients can
dominate device behaviour.

\label{para:thermoelectric_interface} % ---------- %

In thermoelectric materials, interfaces are employed to decouple electrical and thermal transport.

Interface-induced phonon scattering suppresses thermal conductivity while preserving electronic conductivity; enhancing
the thermoelectric figure of merit, $ZT$.

Layered TMDCs and heterostructures with nanostructured interfaces exploit this mechanism to improve thermal management.

\label{para:interfaces_active_modelling_challenge} % ---------- %

Across these examples, it is evident that interfaces are not passive boundaries, but active, tunable regions that
critically determine macroscopic functionality.

Their predictive modelling, through DFT, machine learning potentials, or interface-specific structure search tools like
ARTEMIS, remains a central challenge in modern materials design.

\section{Conventional Prediction Methods} \label{section:conventional_methods}

% ---------- %

Understanding and predicting the behaviour of material interfaces requires a firm foundation in computational methods
that are rooted in quantum and statistical mechanics.

This section introduces the core approaches historically used to model interfaces at the atomic and mesoscale level,
focusing primarily on first-principles methods and classical simulation techniques.

Although each of these methods offers powerful insights, their accuracy, computational cost, and suitability for
interface modelling vary substantially.

% ---------- %

A primary tool in this field is \textit{Density Functional Theory} (DFT), a quantum-mechanical framework capable of
describing electronic structure with high precision.

DFT is well-suited to calculating key interfacial properties such as formation energy, charge redistribution, and band
alignment, making it an essential component of most modern interface studies.

Within this domain, software packages such as \textsc{VASP} have become widely adopted for their efficient
implementations of the Kohn-Sham formalism.

% ---------- %

However, DFT has limits. Computationally, DFT's cubic scaling with system size makes it computationally prohibitive for
large or disordered interface models.

As a result, researchers also turn to \textit{classical methods}, such as \textit{Molecular Dynamics} (MD).

MD enables the simulation of dynamic processes and thermal fluctuations over nanosecond timescales.

It is also adept at capturing thermal effects and atomic rearrangements.

% ---------- %

These conventional approaches provide the groundwork for predicting critical interfacial properties, including adhesion
energy, charge transfer behaviour, and electron or phonon transport.

They remain indispensable despite their limitations, particularly when used in tandem.

This section surveys these methods in turn, addressing their theoretical underpinnings, practical implementations, and
their relative advantages and constraints in the context of interface science.

% ---------- %

A foundational step in many interface modelling workflows is lattice matching between two crystalline materials.

This involves aligning compatible Miller planes to minimise strain and identify viable supercell geometries.

The method introduced by Zur et al.~(1984) laid the groundwork for many subsequent algorithms used in tools such as
\textsc{METADISE} and \textsc{ARTEMIS}, which systematically enumerate low-strain combinations, slips, terminations, and
atomic registries.

These approaches assume idealised surfaces and often neglect more complex processes such as surface reconstructions,
chemical interdiffusion, or local disorder.

% todo: Describe it (lattice matching / geometry-driven interface generation)
% don't just reference.

% ---------- %

DFT remains the most accurate and widely used quantum mechanical method for predicting interfacial properties.

It solves the Kohn--Sham equations to determine the electronic ground state, providing access to total energies, charge
densities, and band edge positions.

% ---------- %

MD simulates the atomic trajectories under classical Newtonian mechanics, governed by interatomic potentials.

Classical MD is well-suited to studying thermal processes such as diffusion, defect migration, and strain relaxation on
picosecond to nanosecond timescales.

% ---------- %

\subsection{Density Functional Theory (DFT)}

% ---------- %

Density Functional Theory (DFT) is a first-principles quantum mechanical method used extensively to model the electronic
structure of condensed matter systems, including molecules, solids, and, critically, material interfaces.

It provides a variational framework in which the ground-state properties of a many-electron system can be obtained from
a functional of the electron density $\rho(\mathbf{r})$, rather than from the many-body wavefunction
$\Psi(\mathbf{r}_1, \dots, \mathbf{r}_N)$.

This simplification, formalised through the Hohenberg-Kohn theorems, underpins the modern implementation of DFT via the
Kohn-Sham approach, where the interacting system is replaced by an auxiliary system of non-interacting electrons subject
to an effective potential that captures Coulomb, exchange, and correlation interactions.

% ---------- %

The total Kohn-Sham energy is given by:

\begin{equation} E_{\mathrm{KS}}[\rho] = T_s[\rho] + \int V_{\mathrm{ext}}(\mathbf{r}) \rho(\mathbf{r}) \, \mathrm{d}\mathbf{r} + \frac{1}{2} \int \int \frac{\rho(\mathbf{r}) \rho(\mathbf{r'})}{|\mathbf{r} - \mathbf{r'}|} \, \mathrm{d}\mathbf{r} \, \mathrm{d}\mathbf{r'} + E_{\mathrm{xc}}[\rho], \end{equation} where $T_s[\rho]$ is the kinetic energy of the non-interacting electrons, $V_{\mathrm{ext}}$ is the external (ionic) potential, and $E_{\mathrm{xc}}[\rho]$ is the exchange-correlation functional, the exact form of which remains unknown.

% ---------- %

Common approximations include the Local Density Approximation (LDA), Generalised Gradient Approximation (GGA), and
hybrid functionals.

While LDA assumes a locally uniform electron gas and often leads to overbinding, GGA incorporates density gradients and
improves accuracy for molecular and inhomogeneous systems.

Hybrid functionals, such as HSE06, incorporate a portion of exact exchange and offer improved bandgap and localisation
predictions at significantly greater computational cost.

% ---------- %

DFT has become foundational in materials modelling due to its ability to predict not only total energies and optimised
geometries, but also derived properties including electronic band structures, charge densities, dielectric responses,
phonon spectra, and interfacial energetics.

Its applications to interface modelling, however, are computationally constrained.

% ---------- %

All DFT simulations in this project are conducted using the Vienna Ab initio Simulation Package (VASP), a widely used
plane-wave-based implementation of Kohn-Sham DFT.

VASP employs the projector-augmented wave (PAW) method to model core-valence interactions, and supports a variety of
exchange-correlation functionals including GGA-PBE and hybrid HSE06.

Structural relaxations are typically performed with conjugate gradient algorithms until residual forces fall below a
chosen threshold (e.g. 0.01 eV/\AA).

% ---------- %

Periodic boundary conditions are employed in all three spatial directions, necessitating the introduction of vacuum
regions for slab models and sufficient supercell dimensions to minimise image-image interactions in interface
simulations.

% ---------- %

The simulation input is controlled via the INCAR, POSCAR, KPOINTS, and POTCAR files, with convergence typically assessed
by k-point density and energy cutoff tests. Pulay and Broyden mixing schemes are used for charge density convergence,
and dipole corrections may be employed for asymmetric slab geometries.

% ---------- %

While VASP provides accurate and robust results for interfacial modelling, it remains computationally intensive.

For extended or complex interfacial systems, this motivates the introduction of machine-learned potentials (MLPs), which
aim to reproduce DFT accuracy.

% ---------- %

\subsection{Other Modelling Methods}

% ---------- %

In addition to DFT, other classes of atomistic modelling, such as MD, has proven valuable for exploring material
interfaces, particularly when temporal evolution or thermal effects are of interest.

% ---------- %

Molecular Dynamics is a time-resolved simulation method in which the trajectories of atoms are computed by numerically
integrating Newton's equations of motion.

Forces on each atom are derived either from empirical interatomic potentials (classical MD) or quantum-mechanical
calculations (ab initio MD, e.g. Car-Parrinello or Born-Oppenheimer approaches).

In interface science, MD is used to study thermal stability, strain relaxation, diffusion, and defect evolution at
finite temperatures.

% ---------- %

One advantage of MD is its capacity to resolve dynamic processes at the atomic level, such as dislocation glide, grain
boundary motion, or interfacial melting.

However, timescales are severely limited: typical simulations span nanoseconds, far shorter than the seconds to minutes
often relevant in experiments.

Classical MD can scale to millions of atoms but sacrifices accuracy; ab initio MD is more accurate but typically limited
to fewer than 1,000 atoms and short timescales due to computational cost.

% ---------- %

For interface modelling, MD is particularly useful when studying thermally activated restructuring, interdiffusion, and
local vibrational behaviour, especially under conditions such as annealing or irradiation.

It also serves as a preliminary stage in workflows involving machine-learned potentials, providing dynamical data for
training.

% ---------- %

MD complements DFT in the modelling of interfaces.

While DFT offers high-fidelity energetics and electronic structure, MD enables atomistic dynamics over nanoseconds.

As such, hybrid approaches are increasingly common, with DFT used to parameterise MD models, or machine learning
techniques introduced to accelerate sampling across scales.

% ---------- %

\subsection{What are the advantages and limitations of DFT and MD?}

% ---------- %

Each computational method traditionally employed in the modelling of interfaces, DFT and MD, offers a distinct balance
between accuracy, computational efficiency, and physical scale.

While often applied in isolation, their complementary regimes suggest the need for hybrid or hierarchical approaches in
the study of real interfaces.

% ---------- %

\textbf{Density Functional Theory (DFT)} is the most widely used \textit{ab initio} method for predicting atomic-scale
properties at interfaces.

Crucially, DFT is capable of modelling materials across the periodic table and is the backbone of modern electronic
structure prediction.

Its limitations include conventional exchange-correlation functionals (e.g. GGA) systematically underestimate band gaps
and cannot capture excited-state dynamics without costly extensions such as GW or time-dependent DFT.

% ---------- %

\textbf{Molecular Dynamics (MD)} offers time-resolved access to thermally activated processes and interfacial dynamics.

Classical MD, using empirical or machine-learned interatomic potentials, enables simulation of thousands to millions of
atoms over nanosecond to microsecond timescales.

However, MD inherits limitations from its force field or potential model: empirical potentials may not transfer well
across bonding environments, while high-fidelity machine-learned potentials require extensive and chemically diverse
training data.

Furthermore, standard MD remains constrained to accessible timescales, making it ill-suited for capturing rare events
such as defect migration, reconstruction, or long-term phase evolution.

% ---------- %

\subsection{How are properties like adhesion energy or electron transport predicted?}

% ---------- %

\textbf{Adhesion energy} is a thermodynamic quantity describing the work required to separate two materials at their
interface.

It is commonly computed from the total energies of the isolated constituents and the relaxed interface configuration.

Within DFT, the adhesion energy $E_\mathrm{adh}$ is given by:

\[ E_\mathrm{adh} = \frac{1}{A} \left( E_\mathrm{slab1} + E_\mathrm{slab2} - E_\mathrm{interface} \right), \]

where $A$ is the interface area, and the energies refer to fully relaxed geometries of the individual slabs and the
interface supercell.

% ---------- %

% todo: add interface energy
% todo: add surface energy
% todo: add strain energy
% todo: add strain: tension, compression, shear
% todo: add formation energy

% ---------- %

\textbf{Electron transport} across interfaces is more nuanced and often requires evaluation of both band alignment and
carrier mobility.

Band alignment can be predicted via several approaches.

The simplest, Anderson's rule aligns vacuum levels to estimate offsets from electron affinities or ionisation
potentials.

However, this often fails to capture interface-specific effects such as charge transfer, built-in fields, or
hybridisation.

More accurate predictions therefore rely on self-consistent DFT calculations that extract the local electrostatic
potential across the heterostructure to determine band offsets directly.

% ---------- %

For example, the offset between conduction or valence band edges across an interface can be obtained by aligning the
macroscopic average electrostatic potentials of the constituent slabs with that of the full interface.

This method accounts for relaxation, dipole formation, and interfacial reconstruction.

In cases where electronic coupling is strong, or the valence states are spatially delocalised across layers (as in
type-II or type-III band alignment), transport characteristics such as charge separation, carrier confinement, or
tunnelling probability can change significantly.

% ---------- %

Charge transport properties such as electrical conductivity, Seebeck coefficient, or mobility may be
approximated from the band dispersion near the Fermi level using Boltzmann transport theory, or more rigorously via
non-equilibrium Green's function (NEGF) methods, especially for low-dimensional or nanoscale interfaces.

% ^^^ note: Intro

\section{Emergence of Machine Learning in Materials Science} \label{section:machine_learning}

% ---------- %

Over the past decade, machine learning (ML) has emerged as a transformative tool in materials modelling, offering a
route to accelerate predictions and expand accessible system sizes beyond those tractable by first-principles methods.

Traditional computational techniques, such as DFT and MD, have established themselves as accurate and rigorous, but they
often scale poorly with system size or configurational complexity.

This limitation becomes particularly acute in the modelling of interfaces, where subtle geometric, chemical, and
electronic variations can dramatically affect material behaviour.

ML offers a complementary paradigm, enabling the construction of surrogate models trained on high-fidelity data, which
can interpolate or even extrapolate to predict key material properties across vast configuration spaces.

% ---------- %

This section introduces the growing role of ML in materials prediction, highlighting its application in interface
modelling and the broader landscape of data-driven approaches.

Subsequent subsections will examine the types of models in use, including classical regressors and deep neural
architectures, alongside the nature of training data, typical input descriptors, and validation strategies.

Special attention will be given to machine-learned interatomic potentials (MLPs), particularly those with demonstrated
success in interfacial contexts such as MACE and ChgNet.

Finally, we reflect on the benefits and ongoing limitations of ML methods in materials science, including issues of
transferability, uncertainty quantification, and integration into automated modelling pipelines.

% ---------- %

\subsection{How has machine learning been introduced into materials prediction?}

% ---------- %

Machine learning (ML) has been introduced into materials prediction as a complementary tool to traditional
first-principles methods, offering improved scalability, predictive flexibility, and the potential for workflow
automation.

Its adoption has been driven by limitations in conventional approaches such as Density Functional Theory (DFT),
particularly in capturing the structural and chemical complexity of real-world systems such as interfaces.

% ---------- %

ML enables efficient exploration of vast configurational landscapes, where brute-force DFT sampling would be infeasible.

This is particularly important in interface modelling, where multiple stacking orders, terminations, and slip
configurations can exist.

ML-driven screening allows researchers to rapidly identify promising candidates for further evaluation.

% ---------- %

Early ML applications focused on surrogate models trained to replicate DFT-derived energies, forces, or properties.

Descriptor-based regression models, using quantities such as coordination numbers, Bader charges, or bond lengths, were
employed to predict interfacial stability or band alignments.

These approaches have evolved toward end-to-end learning frameworks using physically-informed descriptors or graph-based
representations.

% ---------- %

ML is now embedded into iterative prediction--generation--evaluation cycles.

For example, predicted interface favourability informs structure generation (e.g. stacking vectors or surface
terminations), which are then evaluated by DFT or ML surrogates.

The results feed back into model training, supporting a self-refining loop.

Tools such as ARTEMIS benefit from such integration.

% ---------- %

More recent work employs graph neural networks (GNNs), such as MACE or NequIP, which are equivariant to Euclidean
transformations and thus well-suited for atomistic systems.

These models learn directly from atomic positions and species, avoiding the need for handcrafted features while
preserving symmetry constraints essential for interface modelling.

% ---------- %

Ultimately, ML enables a shift from static modelling pipelines to more autonomous workflows.

When coupled with structure-generation protocols and physical constraints.

% ---------- %

\subsection{What types of models and data are being used?}

% ---------- %

Models can be broadly categorised into classical learning algorithms and deep learning approaches.

Classical methods include random forests, gradient boosting, support vector machines (SVMs), and Gaussian process
regression (GPR), the latter being particularly valuable for small datasets due to its uncertainty estimation
capability.

Deep learning architectures include feedforward neural networks and, more notably, graph neural networks (GNNs), which
are well-suited for materials modelling due to their ability to represent variable-sized atomic graphs with complex
bonding topologies.

% ---------- %

The input to these models typically consists of numerical representations of atomic configurations known as descriptors.

Geometric descriptors include bond lengths, coordination numbers, angular distributions, $n$-body distribution
functions, and radial symmetry functions; many of which feature prominently in Behler-Parrinello-type neural networks.

Chemical descriptors such as electronegativity differences and Bader charges capture bonding character and charge
transfer.

Electronic descriptors like partial density of states (PDOS) and local work functions are relevant for predicting band
alignment and interfacial transport properties.

Interface-specific features, such as interfacial roughness or vacuum level offsets, are particularly crucial in 2D|3D
and mixed-dimensional heterostructures.

% ---------- %

To support learning across such heterogeneous and high-dimensional features, symmetry-invariant and scale-independent
representations have become increasingly important.

These ensure that models generalise across materials with different atom counts or lattice sizes.

Dimensionality reduction methods and learned embeddings (e.g. via autoencoders or message passing) are frequently used
to compress structural information into a lower-dimensional latent space optimised for prediction.

% todo: handle note "word soup"...

% ---------- %

The datasets used to train such models vary in size and fidelity.

They often derive from high-throughput DFT calculations available in public repositories such as the Materials Project
or AFLOW.

% ---------- %

\subsection{What are the strengths and limitations of ML in this field?}

% ---------- %

Machine learning (ML) has emerged as a promising strategy for accelerating materials modelling, particularly in the
context of complex interfacial systems.

Its strengths lie chiefly in efficiency, scalability, and adaptability.

ML models, once trained, can rapidly approximate formation energies, surface reactivities, and stability metrics for
thousands of configurations, many of which would be prohibitively expensive to assess using first-principles methods
alone.

This computational efficiency enables exploration of broader configurational and chemical spaces, including systems with
disorder, stacking faults, or metastable structures that pose challenges for traditional simulations.

% ---------- %

Another key advantage is the capacity of ML to capture subtle, non-linear correlations in high-dimensional datasets.

When used in conjunction with tools such as ARTEMIS or RAFFLE, ML can assist not only in predicting energetics, but also
in prioritising candidate geometries during structure generation.

% ---------- %

However, ML approaches also exhibit critical limitations.

A foremost concern is their dependence on the training dataset: models typically interpolate well within the domain of
known examples but generalise poorly to unseen chemistries or extreme geometries.

This limits their reliability when extrapolating to novel material classes or interface types, unless training sets are
curated with care to ensure coverage of relevant chemical and structural motifs.

% ---------- %

Another limitation is interpretability.

While classical models (e.g. those based on density functional theory) yield physically grounded outputs, such as charge
densities or band structures, many ML models operate as statistical black boxes.

Though efforts are underway to integrate physically meaningful descriptors (e.g. n-body distribution functions, local
work function, or Bader charges), ensuring transparency and robustness remains a work in progress.

% ---------- %

Finally, ML alone cannot fully replace physics-based methods.

For systems exhibiting emergent interfacial phases, or electronic reconstructions, high-fidelity methods like DFT or
many-body perturbation theory remain necessary for validation and refinement.

As such, a hybrid strategy is often advocated; using ML to pre-screen candidate structures and flag configurations for
further high-accuracy assessment.

% ---------- %

ML constitutes a valuable addition to the materials modelling toolkit, particularly for screening
and guiding interface prediction.

Yet its performance remains contingent on data quality and physical context, and it is best deployed in tandem with more
rigorous electronic structure methods.

% ---------- %

\subsection{Machine-Learned Interatomic Potentials (MLPs)}

% ---------- %

\subsubsection{MACE}

% ---------- %

The MACE (Many-body Attention and Correlation Equivariant) model represents a recent advancement in equivariant message
passing neural networks (MPNNs), tailored for high-accuracy and scalable force field construction.

Unlike earlier MPNNs which rely primarily on two-body invariant messages, MACE employs many-body interactions and
equivariant tensor contractions to capture angular and chemical complexity with improved expressivity and data
efficiency.

% ---------- %

One notable strength of MACE is its capacity to reach state-of-the-art accuracy on diverse benchmarks, including rMD17,
3BPA, and AcAc, while requiring only two message passing iterations, due to its use of high-body-order messages.

For instance, in extrapolation tasks involving out-of-distribution molecular geometries or high-temperature
trajectories, MACE (particularly with $L=2$) has demonstrated superior performance and robustness relative to competing
models such as NequIP and BOTNet.

% ---------- %

\subsubsection{ChgNet}

% ---------- %

ChgNet is a neural network potential that incorporates both atomic positions and charge distributions into its
prediction framework, thereby explicitly learning charge transfer phenomena.

This is particularly relevant for interfaces where electrostatics, dipole formation, and local charge accumulation
significantly affect electronic properties such as band alignment.


% ---------- %

\subsubsection{Use in Interface Modelling}

% ---------- %

In the context of this project, MLPs like MACE and ChgNet offer distinct but synergistic advantages.

MACE, with its speed, high accuracy, and symmetry-respecting architecture, is well-suited for structure relaxation and
force prediction across multi-material junctions.

ChgNet, on the other hand, may provide more faithful predictions of electronic interface effects arising from local
charge redistribution.

% ---------- %

Therefore, hybrid workflows that use MLPs for broad structural exploration, followed by targeted DFT recalculations for
ambiguous configurations, are increasingly being adopted.

% ---------- %

The eventual goal is to integrate MLPs into an active learning pipeline alongside tools like RAFFLE or ARTEMIS,
facilitating the generation, screening, and optimisation of candidate interface structures at scale.

% todo: address note "should be in intro"

% ---------- %

% ---- where did this come from? ---- %
% todo find if i want to keep this!
Physics-based methods, most notably Density Functional Theory (DFT), provide high accuracy for local bonding
environments and band structure prediction.

However, the computational scaling of standard Kohn-Sham DFT---typically $\mathcal{O}(N^3)$ with respect to system
size---renders large or disordered interface simulations intractable without approximation.

Even when feasible, DFT struggles to capture extended features such as long-range strain relaxation, misfit
dislocations, and spatially diffuse defect states that dominate realistic interface behaviour.

% ---------- %

Machine learning potentials (MLPs), such as those based on graph neural networks or equivariant message passing
architectures, have emerged as promising tools to extend predictive power to larger systems at near-DFT accuracy.

Nonetheless, they too face structural limitations.

MLPs require extensive, representative training datasets, typically derived from DFT calculations, and their performance
deteriorates when applied to configurations involving unseen chemistries, surface reconstructions, or high strain
regimes.

Furthermore, many ML models lack rigorous uncertainty quantification, making it difficult to detect failure modes in
out-of-distribution predictions.

% ---------- %

Both DFT and ML approaches also fall short in capturing emergent interfacial phenomena such as charge transfer,
interface-induced dipoles, or the formation of new interfacial phases.

These effects are often strongly dependent on the local atomic registry and chemical heterogeneity, and cannot be
inferred from the properties of the isolated constituents alone.

Notably, simple models like Anderson's rule have been shown to systematically fail in 2D and 2D|3D heterostructures,
necessitating corrections such as $\Delta E_\Gamma$ and $\Delta E_\mathrm{IF}$ that explicitly depend on the atomic-scale interface
structure and electrostatics.

% ---------- %

While ML and physics-based models have complementary strengths, neither is sufficient in isolation for the reliable
prediction of interfacial behaviour.

Their limitations motivate the development of hybrid workflows, such as embedding physical constraints into ML
architectures, or coupling MLP-based screening with selective DFT refinement.

These approaches aim to achieve a balance between scalability, accuracy, and generalisability---a balance that remains at
the forefront of interface modelling challenges .

% ---^ where did this come from? ^--- %

\section{Limitations of Current Approaches}

% ---------- %

Despite notable advances in both physics-based and machine learning (ML) approaches to materials modelling, several
persistent limitations constrain the predictive fidelity, scalability, and generalisability of current methods,
particularly when applied to realistic interface systems.

This section outlines three critical domains in which these limitations manifest, motivating the hybrid methodologies
explored later in this report:

% ---------- %

\begin{itemize}
  \item \textbf{What are the current bottlenecks in interface prediction?} \\
  Interface prediction remains an inherently multiscale problem, combining combinatorial complexity at the atomic level
  with long-range elastic and electrostatic effects.

  First-principles approaches such as density functional theory (DFT) offer high accuracy but scale cubically with
  system size, making them infeasible for large or disordered systems typical of interfaces.

  Even with machine-learned interatomic potentials (MLPs), exploring configurational space, e.g. over terminations,
  interlayer slips, and misfit reconstructions, requires extensive sampling and often yields flat energy landscapes with
  many metastable minima.

  \item \textbf{Where do both ML and physics-based methods fall short?} \\
  DFT underestimates band gaps and struggles to capture emergent interfacial phenomena such as dipole formation, phase
  reconstruction, or charge transfer without recourse to hybrid functionals or many-body perturbation theory.

  Conversely, ML models, especially those trained on bulk or idealised configurations, exhibit limited transferability
  to chemically complex, strained, or disordered interfaces.

  Without embedded physical constraints or domain-specific descriptors, ML models often fail to resolve interfacial
  states, misfit dislocations, or partial intermixing.

  \item \textbf{Are there data or generalisation issues?} \\
  Interface-specific datasets are significantly less developed than those for bulk crystals, due in part to the vast
  configurational diversity of junctions and the lack of standardised repositories.

  This data sparsity hampers the ability of ML models to generalise across terminations, chemistries, or stacking
  motifs.

  Additionally, many early ML workflows lack symmetry-invariant and scale-independent representations, increasing the
  risk of overfitting to specific atom counts or lattice types.

\end{itemize}

% ---------- %

These limitations underscore the need for hybrid workflows that integrate ML potentials (e.g., MACE, ChgNet) for
large-scale screening, DFT for high-fidelity validation, and structure-generation tools such as ARTEMIS or RAFFLE for
efficient sampling of the configurational landscape.

Such frameworks aim to balance accuracy, efficiency, and transferability in modelling the behaviour of realistic
interfacial systems.

% ---------- %

\subsection{What are the current bottlenecks in interface prediction?}

% ---------- %

The configurational degrees of freedom, ranging from slip vectors, surface terminations, and intermixing patterns to
stacking sequences and misfit accommodation, create a vast structural landscape that resists exhaustive exploration via
conventional methods.

% ---------- %

\paragraph{Configurational Complexity.} At the atomic scale, even between lattice-matched materials, variations in
terminations, registry shifts, and interlayer relaxations produce an exponential number of possible configurations.

The search space rapidly becomes unmanageable, especially in systems requiring large supercells to capture long-range
strain or disorder.

This renders brute-force density functional theory (DFT) calculations infeasible for comprehensive screening.

% ---------- %

\paragraph{Transferability of Machine-Learned Potentials.} Machine-learned interatomic potentials (MLPs), such as MACE
and ChgNet, offer pathways to accelerate structural screening at near-DFT accuracy.

However, their performance is contingent upon the representativeness of the training data.

% ---------- %

\subsection{Where do both ML and physics-based methods fall short?}

% ---------- %

Despite significant progress in both first-principles and machine learning (ML) approaches, key limitations persist in
the predictive modelling of material interfaces.

These limitations arise not only from computational constraints but also from intrinsic mismatches between real
interfacial complexity and the assumptions underpinning current models.

% ---------- %

Physics-based methods, most notably Density Functional Theory (DFT), provide high accuracy for local bonding
environments and band structure prediction.

Even when feasible, DFT struggles to capture extended features such as long-range strain relaxation, misfit
dislocations, and spatially diffuse defect states that dominate realistic interface behaviour.

% ---------- %

MLPs require extensive, representative training datasets, typically derived from DFT calculations, and their performance
deteriorates when applied to configurations involving unseen chemistries, surface reconstructions, or high strain
regimes.

Furthermore, many ML models lack rigorous uncertainty quantification, making it difficult to detect failure modes in
out-of-distribution predictions.

% ---------- %

\subsection{Are there data or generalisation issues?}

% ---------- %

In contrast to bulk crystals, where structured repositories like the Materials Project or AFLOW provide extensive
coverage, interface-specific datasets remain sparse.

This is due to the combinatorial explosion of viable terminations, slip vectors, and orientations across material pairs,
especially in mixed-dimensional systems.

Existing databases tend to overrepresent a narrow set of lattice-matched, technologically prominent interfaces,
introducing bias and hindering the model's capacity to generalise across broader chemical or structural classes.

% ---------- %

Another critical challenge is the lack of standardisation in how interface data is represented.

Unlike bulk materials, which can be described relatively easily by a set of structural features such as lattice
constants and atomic compositions, interfaces require more complex descriptors, including local bond environments,
atomic coordination numbers, and structural motifs that change drastically depending on orientation, termination, and
strain.

As a result, the design of robust and transferable machine learning models for interfaces is further complicated by the
difficulty in developing universally applicable features that capture the intricate nature of interfacial systems.

This problem is compounded by the lack of large, consistent datasets that integrate structural, energetic, and
electronic data, making it harder for models to learn the subtle interactions that govern interface stability and
properties.

% ---------- %

Finally, the sheer diversity of possible interface configurations, especially for systems involving two-dimensional (2D)
materials, complex oxides, or disordered solids, makes high-throughput DFT calculations or other first-principles
methods impractical for exhaustive data collection.

To address these issues, there is a pressing need for hybrid data collection strategies that combine high-fidelity,
physics-based methods with more scalable, data-driven approaches, such as machine learning potential models that can
rapidly generate large datasets of interfacial configurations.

Additionally, targeted experimental datasets that validate theoretical predictions are essential for creating reliable
and comprehensive interfaces databases.

% ---------- %

Generalisation error is further amplified by suboptimal representation of interface-specific physics.

Many descriptor schemes, such as SOAP vectors or radial symmetry functions, are optimised for bulk periodicity and fail
to capture interfacial asymmetry, vacuum discontinuities, or broken symmetry.

Representations that do not enforce scale-invariance or rotational equivariance can also hinder model transferability to
supercells or faceted configurations.

% ---------- %

Finally, inconsistencies in training labels present an often-overlooked source of predictive error.

Energies obtained from high-throughput DFT pipelines can vary depending on pseudopotential choice, relaxation
tolerances, or convergence thresholds.

In interfacial contexts, these uncertainties are magnified by metastability and complex relaxation pathways.

Some datasets rely on static, unrelaxed geometries or surrogate models, further undermining the reliability of ground
truth energetics.

% ---------- %

While ML offers substantial speed-ups over DFT, its utility in interface prediction is presently constrained by data
bottlenecks, representational mismatch, and transferability failures.

Addressing these limitations will require a multi-pronged strategy involving broader and more diverse training datasets,
improved descriptors attuned to interfacial physics, and integration of uncertainty quantification into ML workflows.

Hybrid schemes, where MLPs perform large-scale screening followed by selective DFT refinement, may offer a viable path
forward.

\section{State-of-the-Art in Interface Prediction} \label{section:state_of_the_art}

% ---------- %

Recent years have seen rapid developments in the modelling of interface structures, driven both by methodological
advances and by the growing technological importance of heterostructured materials.

Interface prediction now spans a range of computational approaches, from first-principles calculations and
high-throughput screening to data-driven and machine learning (ML) methods.

This section provides an overview of the current state-of-the-art, establishing the context for the subsequent
discussion of recent studies (Section 2.5.1), the positioning of this project's methodology (Section 2.5.2), and
anticipated emerging directions (Section 2.5.3).

% ---------- %

A key enabler of recent progress has been the development of automated interface generation tools such as ARTEMIS and
InterMatch, which systematically enumerate possible lattice matches, terminations, and stacking configurations between
two crystal structures.

When integrated with rapid energy evaluation, either via density functional theory (DFT) or machine-learned potentials
(MLPs), these tools allow the configurational space of interface structures to be sampled more broadly and with greater
fidelity.

In parallel, studies have increasingly sought to move beyond idealised, coherent interfaces to include semi-coherent,
faceted, and amorphous boundaries, reflecting the complexity of real materials under fabrication conditions.

% ---------- %

Efforts have also been made to improve the prediction of interfacial electronic properties, such as band alignment.

Where traditional rules like Anderson's rule fail, corrective models based on physically grounded terms, such as
$\Delta E_{\Gamma}$ and $\Delta E_{\mathrm{IF}}$, have been proposed and validated across a range of 2D materials.

Such corrections are essential for capturing charge transfer and hybridisation effects at interfaces, and form an
important bridge between local structure and electronic response.

% ---------- %

This chapter aims to synthesise these various advances, situating the present work within an evolving landscape of
interface modelling.

Particular attention will be given to scalable workflows that combine structural generation, energetic evaluation, and
property prediction, setting the stage for a more autonomous and predictive framework for materials interface discovery.

% ---------- %

Recent years have witnessed a pronounced expansion in methodologies for predicting interface structures, driven by the
convergence of first-principles modelling, high-throughput screening, and data-driven inference.

% ---------- %

This evolution reflects a shift from idealised, geometry-constrained models towards scalable frameworks capable of
capturing configurational complexity, long-range strain, and emergent phenomena at interfaces.

% ---------- %

Hybrid methodologies now dominate the landscape of ab initio interface prediction. Notable among them is the
\textsc{ARTEMIS} tool developed by Taylor et al., which automates the generation of candidate interface structures by
identifying lattice-matched slabs, multiple surface terminations, and lateral shifts across a wide range of Miller
planes.

% ---------- %

Shortly after, Taylor et al.\ introduced the \textsc{RAFFLE}, with Pitfield et al.\ introducing methodology, combining
empirical void-filling algorithms with iterative statistical descriptors to model metastable and low-symmetry
reconstructions.

% ---------- %

Constrained variants of the Minima-Hopping Method (MHM) allow for controlled exploration of grain boundary
reconstructions by varying atomic densities and introducing geometric constraints.

These approaches yield more realistic reconstructions than fixed-stoichiometry DFT or classical potentials, especially
in systems like silicon where interface states are sensitive to atomic arrangement and misfit accommodation.

% ---------- %

Machine learning (ML) is increasingly embedded into the prediction workflow. Gerber et al.\ developed the
\textsc{InterMatch} framework, which leverages high-throughput DFT databases to train ML models that guide lattice
matching and stacking predictions in 2D systems. At the electronic level, Davies et al.\ proposed physically motivated
corrections to band alignment, namely $\Delta E_\Gamma$ and $\Delta E_{\mathrm{IF}}$, to extend Anderson's rule and
account for hybridisation and interface dipoles in 2D heterostructures.

% ---------- %

Despite these advances, several limitations persist.

Data scarcity for interfacial systems hampers generalisation, and most ML models are trained on bulk-like environments.

Existing descriptors may fail to capture the broken symmetry, local dipoles, and charge redistribution unique to
interfaces.

Furthermore, many workflows truncate the configurational space prematurely or rely on geometric heuristics alone.

% ---------- %

There is increasing recognition that hybrid workflows, combining ML surrogates with tools like \textsc{ARTEMIS} and
\textsc{RAFFLE}, and validating critical configurations using DFT, offer a scalable route forward.

This multiresolution strategy enables both efficient exploration and physically grounded evaluation of interfacial
stability across diverse systems.

% ---------- %

\subsection{How does my proposed approach compare and build upon it?}

% ---------- %

Contemporary approaches to interface structure prediction combine symmetry-based lattice matching, high-throughput
screening, and density functional theory (DFT) relaxation.

Tools such as \textsc{ARTEMIS} enable the systematic generation of interface candidates by identifying lattice
commensurabilities and constructing multiple terminations and stacking arrangements from parent surfaces.

However, the energetic evaluation of these candidates remains a bottleneck: full relaxation and total energy comparison
using DFT is prohibitively expensive for systems exceeding a few hundred atoms.

% ---------- %

To mitigate this, recent advances have introduced surrogate strategies that bypass exhaustive DFT sampling.

InterMatch, a high-throughput framework proposed by Gerber \textit{et al.}, guides stacking predictions using
pre-computed DFT databases and electrostatic heuristics.

Yet such frameworks often lack spatial resolution or generalisability across chemically diverse or defective systems.

% ---------- %

This project builds directly upon these developments, advancing them in three ways:

% ---------- %

\begin{enumerate}

   \item \textbf{Descriptor-led generalisation.} Inspired by the success of band-alignment corrections in TMDC
   heterostructures, this work seeks to identify whether interfacial stability can be inferred from properties
   computable on isolated surfaces; such as cleavage ene.g. surface dipole, or projected local density of states near
   the Fermi level.

   If validated, such descriptors would permit extrapolation to unseen terminations or chemistries without requiring
   full interface models.

   \item \textbf{Feedback into structure-generation pipelines.} In the longer term, this project aims to integrate
   ML-derived stability metrics into tools like \textsc{ARTEMIS} or \textsc{RAFFLE}, allowing early pruning of
   unfavourable configurations before computationally costly evaluation.

   This hybrid generation--relaxation loop reflects a broader trend towards data-assisted automation in interfacial
   materials modelling.

\end{enumerate}

% ---------- %

The proposed framework thus combines multi-resolution structure prediction with predictive modelling rooted in
physically meaningful descriptors.

Rather than replacing existing DFT-centric approaches, it seeks to augment them; by embedding surrogate models within
the structure-generation workflow and using them to navigate large configurational spaces.

This enables a shift from retrospective evaluation to prospective design of realistic, energetically favourable
interfaces.

% ---------- %

\subsection{Emerging trends}

% ---------- %

Recent advances in interface prediction have increasingly centred on integrating data-driven approaches with atomistic
modelling frameworks.

Machine-learned potentials (MLPs) such as MACE and E(3)-equivariant networks like NequIP are now enabling the efficient
evaluation of interfacial configurations across large configurational spaces.

This shift permits relaxation of structures at near-DFT accuracy, allowing researchers to explore interface separations,
registry shifts, and metastable motifs that would be prohibitively expensive using standard ab initio methods.

% ---------- %

Collectively, these approaches suggest a field increasingly focused on scalable, automatable, and physically
interpretable modelling of interfaces.

The combination of machine learning with established electronic structure theory is opening a path toward predictive
frameworks that are sensitive to real-world synthesis constraints and device contexts.